\setchapterpreamble[u]{\margintoc}
\chapter{Eigenvalues and Eigenvectors}
\labch{eigenvalues}

\sources{Strang, \emph{Introduction to Linear Algebra}, 6th ed.\ \cite{strang2023ila},
§6.1--6.5; MIT 18.06 \cite{ocw1806}, Lectures~21--22 and 26--28.}

Everything so far treated \(A\) as a machine for solving. This chapter asks a
different question: are there directions in which \(A\) does nothing more
interesting than stretching? The answer reorganises the matrix around those
directions and makes repeated application easy to describe.

\section{The defining equation}
\labsec{eigen-definition}

\begin{definition}
A non-zero vector \(\vect{x}\) is an \emph{eigenvector} of \(A\in\R^{n\times n}\)
with \emph{eigenvalue} \(\lambda\) when
\[
	A\vect{x}=\lambda\vect{x}.
\]
\end{definition}

Rewriting as \((A-\lambda I)\vect{x}=\vect{0}\) shows that \(\lambda\) is an
eigenvalue exactly when \(A-\lambda I\) is singular, since a non-zero vector must
sit in its nullspace. By \refch{determinants} that is the condition
\[
	\det(A-\lambda I)=0,
\]
the \emph{characteristic equation}. Its left-hand side is a polynomial of degree
\(n\) in \(\lambda\).

\marginnote{Note the order of work: first find the \(\lambda\) that make the
matrix singular, then find the nullspace for each. Eigenvectors are never found
first.}

\begin{example}
For \(A=\begin{bmatrix}2&1\\1&2\end{bmatrix}\),
\[
	\det\begin{bmatrix}2-\lambda&1\\1&2-\lambda\end{bmatrix}
	=(2-\lambda)^{2}-1=\lambda^{2}-4\lambda+3,
\]
with roots \(\lambda_1=3\) and \(\lambda_2=1\). Solving
\((A-3I)\vect{x}=\vect{0}\) gives \(\vect{x}_1=[1,1]\tp\), and
\((A-I)\vect{x}=\vect{0}\) gives \(\vect{x}_2=[1,-1]\tp\). The two eigenvectors
are orthogonal, which \refsec{spectral} explains.
\end{example}

\section{Trace and determinant}
\labsec{trace-det}

Two shortcuts follow from comparing coefficients in the characteristic
polynomial.

\begin{proposition}
The eigenvalues of \(A\), counted with multiplicity, satisfy
\[
	\sum_{i}\lambda_i=\operatorname{tr}A=\sum_i a_{ii},
	\qquad
	\prod_i\lambda_i=\det A .
\]
\end{proposition}

For the example, \(3+1=4=2+2\) and \(3\cdot1=3=4-1\). In two dimensions these two
equations determine the eigenvalues outright, which is often faster than forming
the polynomial.

\begin{remark}
Singularity and a zero eigenvalue are the same statement, since the product of
the eigenvalues is the determinant. The nullspace of \(A\) is precisely the
eigenspace for \(\lambda=0\).
\end{remark}

\section{Diagonalisation and powers}
\labsec{diagonalisation}

Suppose \(A\) has \(n\) independent eigenvectors. Placing them in the columns of
\(X\) and the eigenvalues on the diagonal of \(\Lambda\),
\[
	AX=X\Lambda
	\quad\Longrightarrow\quad
	A=X\Lambda X^{-1}.
\]

The payoff is that powers telescope:
\[
	A^{k}=X\Lambda^{k}X^{-1},
\]
because the inner \(X^{-1}X\) factors cancel. Raising a diagonal matrix to a
power is done entry by entry, so the long-run behaviour of \(A^{k}\) is decided
entirely by the sizes \(|\lambda_i|\).

\begin{example}
Let \(A=\begin{bmatrix}0.8&0.3\\0.2&0.7\end{bmatrix}\), whose columns sum to~1.
Its trace is \(1.5\) and its determinant is \(0.5\), so the eigenvalues solve
\(\lambda^{2}-1.5\lambda+0.5=0\), giving \(\lambda_1=1\) and \(\lambda_2=0.5\).
The eigenvector for \(\lambda_1=1\) is \([3,2]\tp\); normalising it to sum to~1
gives \([0.6,0.4]\tp\). Since \(0.5^{k}\to0\), every starting distribution
converges to that vector.
\end{example}

\marginnote{This is the entire theory of Markov chains in miniature: an
eigenvalue of exactly~1 supplies the steady state, and the second largest
modulus sets the rate of approach.}

Diagonalisation can fail. The matrix \(\begin{bmatrix}0&1\\0&0\end{bmatrix}\)
has \(\lambda=0\) twice but only a one-dimensional eigenspace, so no basis of
eigenvectors exists. Distinct eigenvalues guarantee independence and therefore
diagonalisability; repeated eigenvalues may or may not.

\section{Symmetric matrices}
\labsec{spectral}

Symmetric matrices are the well-behaved case, and they are the case that occurs
in practice --- covariance matrices, Gram matrices \(A\tp A\), and Hessians are
all symmetric.

\begin{theorem}[Spectral theorem]
If \(A\tp=A\) then all eigenvalues of \(A\) are real, eigenvectors for distinct
eigenvalues are orthogonal, and \(A\) can be written
\[
	A=Q\Lambda Q\tp
\]
with \(Q\) orthogonal and \(\Lambda\) real diagonal.
\end{theorem}

\begin{proof}[Orthogonality]
Let \(A\vect{x}=\lambda\vect{x}\) and \(A\vect{y}=\mu\vect{y}\) with
\(\lambda\neq\mu\). Then
\[
	\lambda\,\vect{y}\tp\vect{x}=\vect{y}\tp A\vect{x}
	=(A\vect{y})\tp\vect{x}=\mu\,\vect{y}\tp\vect{x},
\]
using symmetry in the middle step. Hence
\((\lambda-\mu)\vect{y}\tp\vect{x}=0\), and since \(\lambda\neq\mu\) the inner
product vanishes.
\end{proof}

Because \(Q^{-1}=Q\tp\), the factorisation writes \(A\) as a sum of orthogonal
rank-one pieces,
\[
	A=\sum_{i=1}^{n}\lambda_i\,\vect{q}_i\vect{q}_i\tp ,
\]
which is the column-times-row reading of \refsec{four-ways}, applied here to
the basis of eigenvectors.

\section{Positive definite matrices}
\labsec{positive-definite}

\begin{definition}
A symmetric \(A\) is \emph{positive definite} when \(\vect{x}\tp A\vect{x}>0\)
for every \(\vect{x}\neq\vect{0}\), and \emph{positive semidefinite} when the
inequality is weak.
\end{definition}

\begin{proposition}
For a symmetric matrix \(A\) the following are equivalent: every eigenvalue is
positive; every pivot is positive; \(\vect{x}\tp A\vect{x}>0\) for every
\(\vect{x}\neq\vect{0}\); and \(A=B\tp B\) for some \(B\) with independent
columns.
\end{proposition}

\begin{example}
\(A=\begin{bmatrix}2&1\\1&2\end{bmatrix}\) has eigenvalues \(3\) and \(1\),
pivots \(2\) and \(\tfrac32\), and
\(\vect{x}\tp A\vect{x}=2x_1^{2}+2x_1x_2+2x_2^{2}=(x_1+x_2)^{2}+x_1^{2}+x_2^{2}\),
positive unless both components vanish. All three tests agree.
\end{example}

\marginnote{Positive definiteness is the matrix version of ``strictly convex''.
\refch{gradient-methods} uses exactly this: a quadratic with a positive definite
Hessian has a single minimum, and the ratio of largest to smallest eigenvalue
predicts how slowly gradient descent will find it.}

\section{Similar matrices}
\labsec{similar}

Matrices \(A\) and \(M^{-1}AM\) are called \emph{similar}. They have the same
characteristic polynomial and hence the same eigenvalues, because
\[
	\det(M^{-1}AM-\lambda I)=\det\bigl(M^{-1}(A-\lambda I)M\bigr)=\det(A-\lambda I).
\]
Similarity is a change of basis, so the message is that eigenvalues belong to
the underlying transformation rather than to the particular array of numbers
used to record it --- a point taken up properly in
\refch{linear-transformations}.
